\PassOptionsToPackage{table,xcdraw}{xcolor}
\documentclass[11pt, a4paper]{article}

\usepackage[T1]{fontenc}
\usepackage[utf8]{inputenc}
\usepackage{tgpagella}
\usepackage[spanish, es-noshorthands]{babel}
\usepackage{amsmath, amssymb}
\usepackage{booktabs}
\usepackage{tabularx}
\usepackage{multicol}
\usepackage[most]{tcolorbox}
\usepackage[american]{circuitikz}
\usepackage[margin=2cm]{geometry}

\definecolor{ucbblue}{RGB}{0, 51, 102}

% Instructor note:
% The two-battery option is intentionally supported so students
% can complete practical work outside the university laboratory.
% Do not penalize battery-powered implementation if actual rails
% are documented and device limits are respected.

\begin{document}

\begin{center}
    \Large\textbf{\textcolor{ucbblue}{Suplemento de Guía: Opciones de Fuente de Alimentación}}
\end{center}

\begin{tcolorbox}[colback=yellow!5, colframe=ucbblue, title=\textbf{Opciones de Fuente de Alimentación}]
    Este laboratorio puede realizarse utilizando una \textbf{fuente dual de banco} o \textbf{dos baterías de $9\,\text{V}$} conectadas en serie. Ambas son opciones de implementación legítimas. En ambos casos, el circuito debe proveer un riel positivo ($V_+$), una referencia común ($0\,\text{V}$) y un riel negativo ($V_-$).
\end{tcolorbox}

\section*{1. Configuración de la Fuente de Alimentación}

\subsection*{Opción A: Fuente Dual de Banco}
Si dispone de equipo de laboratorio, configure la fuente dual para entregar nominalmente:
\begin{equation*}
    \boxed{+10\,\text{V}, \quad 0\,\text{V}, \quad -10\,\text{V}} \quad \text{o} \quad \boxed{+9\,\text{V}, \quad 0\,\text{V}, \quad -9\,\text{V}} \text{}
\end{equation*}
(Siempre verifique primero que el Op-Amp soporte estos voltajes según su datasheet).

\subsection*{Opción B: Dos Baterías de 9 V en Serie (Uso Portátil/Casa)}
\begin{minipage}{0.5\textwidth}
    Utilice dos baterías independientes de $9\,\text{V}$. El punto de unión entre ambas forma la referencia del circuito ($0\,\text{V}$).
    
    \vspace{0.3cm}
    \textbf{Conexión:}
    \begin{itemize}
        \item Batería 1 (+): Conectar a riel $V_+$.
        \item Batería 1 (-): Conectar a la unión central ($0\,\text{V}$).
        \item Batería 2 (+): Conectar a la unión central ($0\,\text{V}$).
        \item Batería 2 (-): Conectar a riel $V_-$.
    \end{itemize}
    \begin{tcolorbox}[colback=red!5, colframe=red!70!black]
        \textbf{Nota Crítica:} El punto medio es la referencia $0\,\text{V}$. \textbf{NO} trate el terminal negativo de la batería inferior como "tierra".
    \end{tcolorbox}
\end{minipage}
\hfill
\begin{minipage}{0.45\textwidth}
    \begin{center}
        \begin{circuitikz}[scale=0.9, transform shape, thick]
            % Bateria 1 (Superior)
            \draw (0,3) to[battery2, l=$9\,\text{V}$ (Bat 1)] (0,1);
            \draw (0,3) to[short, *-o] (2,3) node[right]{$V_+ \approx +9\,\text{V}$};
            
            % Union central y tierra
            \draw (-1,1) to (0,1) to [short, *-o] (2,1) node[right]{$0\,\text{V}$ (Midpoint)};
            \draw (-1,1) node[ground]{};
            
            % Bateria 2 (Inferior)
            \draw (0,1) to[battery2, l=$9\,\text{V}$ (Bat 2)] (0,-1);
            \draw (0,-1) to[short, *-o] (2,-1) node[right]{$V_- \approx -9\,\text{V}$};
        \end{circuitikz}
    \end{center}
\end{minipage}

\section*{2. Realidad de las Baterías y Medición Obligatoria}
Una batería nueva puede medir más de $9\,\text{V}$, y su voltaje caerá con el uso. Es común que ambas baterías tengan voltajes distintos, resultando en rieles \textbf{asimétricos} (ej. $V_+ = +9.2\,\text{V}$, $V_- = -8.7\,\text{V}$). \textbf{Debe medir y registrar los voltajes reales; no asuma valores nominales exactos.} En LTSpice, si desea alta fidelidad, actualice las fuentes con estos valores medidos.

\section*{3. Reglas de Compatibilidad y Seguridad}
\begin{itemize}
    \item \textbf{Desacoplo Obligatorio:} Las baterías \textbf{no} eliminan la necesidad de usar capacitores de desacoplo ($100\,\text{nF}$) cerca de los pines del Op-Amp hacia $0\,\text{V}$.
    \item \textbf{Verificación del Datasheet:} No todos los Op-Amps operan igual a $\pm9\,\text{V}$. Identifique la familia exacta (LM741, UA741CN, LM358N, TL074), verifique el rango de alimentación, el \textit{output swing} y el \textbf{pinout}. ¡No son compatibles pin-a-pin!
    \item \textbf{Seguridad Casera:} Desconecte las baterías cuando no use el circuito, evite cortocircuitos, y verifique la polaridad antes de insertar el CI.
\end{itemize}

\section*{4. Adaptación Dinámica Crítica para Lab 4 (Slew Rate)}
Para la prueba de gran señal, el objetivo es $V_p = 4\,\text{V}$ ($8\,\text{V}_{pp}$). Siga esta lógica de decisión obligatoria para evitar confundir saturación con Slew Rate:
\begin{enumerate}
    \setlength{\itemsep}{0pt}
    \item Mida los rieles reales (banco o baterías).
    \item A baja frecuencia ($1\,\text{kHz}$), intente ajustar la salida a $V_{out} \approx 8\,\text{V}_{pp}$.
    \item \textbf{Revise recortes (Clipping):} ¿La onda se aplasta arriba o abajo? Si los rieles (ej. baterías descargadas) o el \textit{output swing} del Op-Amp no permiten $8\,\text{V}_{pp}$ limpios, \textbf{reduzca la amplitud de salida}.
    \item Si debió reducir la amplitud, registre el nuevo $V_p$ \textbf{real} alcanzado sin distorsión y \textbf{recalcule} sus predicciones teóricas:
    \begin{equation*}
        SR_{\mathrm{req}} = 2\pi f V_{p,\mathrm{nuevo}} \quad \text{y} \quad f_{\max,SR} = \frac{SR_{\mathrm{datasheet}}}{2\pi V_{p,\mathrm{nuevo}}} \text{}
    \end{equation*}
    \item \textbf{Solo entonces} comience a subir la frecuencia para observar la verdadera distorsión por Slew Rate (triangularización de la onda).
\end{enumerate}

\section*{5. Tabla de Registro de Fuente (Para Informe IEEE)}
\textit{Añada esta información a su informe IEEE. Especifique qué fuente utilizó en su sección de metodología.}

\begin{table}[h!]
    \centering
    \footnotesize
    \renewcommand{\arraystretch}{1.5}
    \begin{tabularx}{0.8\textwidth}{|>{\hsize=1.2\hsize}X|>{\hsize=0.8\hsize}X|}
    \hline
    \textbf{Variable de Alimentación} & \textbf{Valor Medido Real} \\ \hline
    Fuente utilizada & [Banco / $2 \times$ Bat. $9\,\text{V}$] \\ \hline
    Voltaje Batería 1 ($V_{B1}$) [Si aplica] & \rule{3cm}{0.4pt} \\ \hline
    Voltaje Batería 2 ($V_{B2}$) [Si aplica] & \rule{3cm}{0.4pt} \\ \hline
    Riel Positivo medido ($V_+$) & \rule{3cm}{0.4pt} \\ \hline
    Riel Negativo medido ($V_-$) & \rule{3cm}{0.4pt} \\ \hline
    Alimentación Total ($V_{TOTAL} = V_+ - V_-$) & \rule{3cm}{0.4pt} \\ \hline
    \end{tabularx}
\end{table}

\end{document}
